\section{Introduction}
\label{sec:introduction}

Earthquake-triggered landslides can affect large, inaccessible areas within minutes. A useful response map must arrive before field reconnaissance and before cloud-free optical acquisitions are guaranteed. SAR is suited to this interval because it operates through cloud cover and independent of daylight \cite{Mondini2021,Handwerger2022,Amatya2023}. Yet, the side-looking geometry changes which slope evidence is observed, making the most difficult cases more than ordinary semantic ambiguity.

Foreshortening, layover, and shadow arise from the relation between terrain and radar look direction, particularly in mountainous areas \cite{Mondini2021,Shi2024}. Backscatter and speckle further make local landslide signatures view dependent \cite{Santangelo2022}. Ascending and descending passes thus provide complementary, rather than interchangeable, observations. Combining both geometries improved SAR-change products for rapid landslide density mapping \cite{Handwerger2022}. Yet the complementary post-event pass may not have arrived when the first response map is required. A system that always consumes two post-event orbits trades away the deployment condition that motivates SAR in the first place.

We consider dual-orbit training with single-orbit deployment. This is related to learning with privileged information \cite{LopezPaz2016} and missing-modality distillation \cite{Li2022,Liu2024}, which generally transfer features, logits, or representations. The unanswered question is different: how can an unavailable complementary observation teach a model to better exploit information that remains available? Reconstructing counter-orbit features is not the natural objective because some view-dependent evidence is absent from the target measurement.

Deformable operators provide a useful mechanism for separating this question into search policy and searched values. Deformable convolution and attention learn a sparse set of sampling locations and aggregation weights \cite{Dai2017,Zhu2021}; remote-sensing segmentation has already benefited from adaptive spatial aggregation \cite{Zuo2022,Dai2023}. We use complementary context only to determine \emph{where} and \emph{how strongly} a teacher should search. The searched values, and therefore the decoder input, always come from the target orbit. This excludes a dual-orbit value-fusion bypass that a single-orbit student could not execute.

We propose counter-guided target evidence search for post-event SAR landslide segmentation. The Counter-Guided Search Teacher sees target and counter orbits but uses the counter orbit only to generate a target-side deformable search policy. Its decoder still reads target values only. The Target-Only Search Student learns the policy and the evidence retrieved by that policy from its own target memory. In short, we do not transfer unavailable counter-orbit information; we transfer how to search useful evidence from the target orbit. The contributions are:

\begin{itemize}
    \item A training-dual-orbit and deployment-single-orbit formulation in which the counter orbit is privileged guidance, not an inference input.
    \item A counter-guided target-evidence search mechanism that separates policy generation from searched values, keeping the complete value path target-only.
    \item A target-only student that distills search behavior through policy and shared-memory evidence alignment, evaluated on geometry-valid and full domains with search and loss-component diagnostics.
\end{itemize}
